\documentclass[11pt, a4paper]{article}

% --- PACKAGES ---
\usepackage[margin=1in]{geometry}
\usepackage{amsmath, amssymb} % For math symbols
\usepackage{graphicx}
\usepackage[colorlinks=true, urlcolor=blue, linkcolor=blue]{hyperref}
\usepackage{fancyhdr}
\usepackage{booktabs}   % For professional tables
\usepackage{multicol}   % For multi-column layout

% --- DOCUMENT INFORMATION ---
\title{Lecture 2: The Engine of Learning — Minimizing Loss with Stochastic Gradient Descent}
\author{A Course Outline and Motivation}
\date{\today}

% --- HEADER/FOOTER ---
\pagestyle{fancy}
\fancyhf{}
\rhead{Lecture 2: Optimization with SGD}
\lhead{\thepage}

\begin{document}

\maketitle
\thispagestyle{empty}
\tableofcontents
\newpage

% =================================================================
\section{Recap and Our Next Question}
% =================================================================
In our last lecture, we saw that a deep neural network is a powerful, flexible function built from simple, stacked layers. We represented all the tunable ``knobs" of this function—its weights and biases—with a single giant vector, which we'll call $\mathbf{w}$.

Our goal is to find the best possible settings for $\mathbf{w}$ to solve a task, like classifying images. But this raises a crucial question:
\begin{center}
\Large\textbf{How do we know what ``best" means, and how do we find it?}
\end{center}
Today, we'll answer this by introducing the two core components of the learning process: the \textbf{loss function}, which measures how badly our network is performing, and \textbf{Stochastic Gradient Descent (SGD)}, the algorithm we use to systematically improve it.

% =================================================================
\section{Measuring Performance: The Loss Function}
% =================================================================
To improve our network, we first need a way to quantify its error. This is the job of the \textbf{loss function}, $L(\mathbf{w})$. It takes our network's current weights $\mathbf{w}$ and tells us a single number representing the total error across all our training examples.

Most loss functions in machine learning are a big sum (or average) over the individual errors for each training sample:
\begin{equation}
\min_{\mathbf{w}} L(\mathbf{w}) = \min_{\mathbf{w}} \frac{1}{N} \sum_{i=1}^{N} \ell_i(\mathbf{w})
\end{equation}
Where:
\begin{itemize}
    \item $\mathbf{w} \in \mathbb{R}^d$ are the model parameters we want to find (our weights and biases).
    \item $N$ is the number of training samples (often in the millions).
    \item $d$ is the number of parameters (can be in the billions for large models!).
    \item $\ell_i(\mathbf{w})$ is the loss on the $i$-th individual training sample. It measures the error our network with weights $\mathbf{w}$ makes on that one example.
\end{itemize}
The entire goal of ``training" a neural network is to solve this optimization problem: find the weights $\mathbf{w}$ that make the total loss $L(\mathbf{w})$ as small as possible.

% =================================================================
\section{How to Minimize Loss? The Path of Steepest Descent}
% =================================================================
Imagine you're standing on a hilly landscape, and your goal is to get to the lowest point. The most intuitive strategy is to look around, find the direction that goes downhill most steeply, and take a small step in that direction. You repeat this process over and over until you reach the bottom.

This is exactly the idea behind \textbf{Gradient Descent}. The ``direction of steepest descent" is given by the negative gradient of the loss function, $-\nabla L(\mathbf{w})$. The algorithm is beautifully simple:
\begin{equation}
\mathbf{w}_{\text{new}} = \mathbf{w}_{\text{old}} - \alpha \nabla L(\mathbf{w}_{\text{old}})
\end{equation}
Here, $\alpha$ is the \textbf{learning rate}—it controls how big of a step we take.

\subsection{The Computational Bottleneck}
This seems like a perfect solution! However, there's a huge practical problem. To calculate the true gradient, we have to consider every single data point:
\begin{equation}
\nabla L(\mathbf{w}) = \frac{1}{N} \sum_{i=1}^{N} \nabla \ell_i(\mathbf{w})
\end{equation}
For a modern dataset with millions of samples ($N$ is large), computing this sum just to take \emph{one single step} is incredibly slow. It could take hours! If it takes thousands of steps to train our model, this approach becomes completely infeasible. We need a faster way.

% =================================================================
\section{The Solution: Stochastic Gradient Descent (SGD)}
% =================================================================
The idea of SGD is simple: instead of calculating the exact, expensive gradient using all $N$ samples, what if we approximate it using just \textbf{one randomly chosen sample} at each step?

The SGD update rule is:
\begin{equation}
\mathbf{w}_{\text{new}} = \mathbf{w}_{\text{old}} - \alpha \nabla \ell_{i_k}(\mathbf{w}_{\text{old}})
\end{equation}
where $i_k$ is an index picked randomly from our $N$ samples at step $k$. Each step is now roughly $N$ times faster!

\subsection{Why Does This Even Work? The Unbiased Estimate}
This might seem reckless. How can a gradient from a single, random sample possibly guide us in the right direction? The mathematical magic lies in the fact that this ``stochastic gradient" is an \textbf{unbiased estimator} of the true gradient.

\begin{equation}
\mathbb{E}_{i_k}[\nabla \ell_{i_k}(\mathbf{w})] = \nabla L(\mathbf{w})
\end{equation}
In plain English: while any single stochastic step might be a bit ``off," on average, these steps point in the same direction as the true, full-batch gradient. We've traded a slow, precise step for many fast, noisy-but-correct-on-average steps.

\subsection{The Two Zones of SGD: Geometric Intuition}
The noisy nature of SGD gives it a characteristic two-phase behavior:

\begin{description}
    \item[The ``Far-Out" Zone:] When our model is very wrong (i.e., far from the optimal weights), almost any random sample we pick will tell us to move in roughly the same, correct direction. In this phase, SGD makes rapid, consistent progress, much like a speedboat heading toward a distant shore.
    \item[The ``Region of Confusion":] As our model gets closer to the optimal weights, the situation changes. Now, different random samples might give conflicting directions. One sample might say ``go left," another ``go right." The algorithm starts to ``bounce around" the minimum.
\end{description}

This noisy bouncing isn't necessarily a bad thing! Many researchers believe this stochasticity helps the model avoid getting stuck in poor local minima and can lead to solutions that \textbf{generalize better} to new, unseen data. It prevents the model from ``memorizing" the training set too perfectly.

\subsection{Practical SGD: Using Minibatches}
In practice, we rarely use a single sample. Instead, we use a small \textbf{minibatch} of samples (e.g., 32, 64, or 256) at each step:
\begin{equation}
\mathbf{w}_{\text{new}} = \mathbf{w}_{\text{old}} - \alpha \left( \frac{1}{B} \sum_{j \in \text{Batch}_k} \nabla \ell_j(\mathbf{w}_{\text{old}}) \right)
\end{equation}
This is the best of both worlds:
\begin{enumerate}
    \item \textbf{Reduces Noise:} Averaging over a small batch gives a more stable gradient estimate than a single sample.
    \item \textbf{Computationally Efficient:} Modern hardware like GPUs are designed for parallel computation and are much faster at processing a small batch of 32 examples at once than processing 32 examples one-by-one.
\end{enumerate}

% =================================================================
\section{Our Next Steps: The Course Ahead}
% =================================================================
We've just uncovered the core engine of deep learning. Here’s what we'll explore next:

\begin{description}
    \item[How do we compute the gradients $\nabla \ell_i(\mathbf{w})$?] This is the final missing piece. In our next lecture, we'll learn about \textbf{backpropagation}, the remarkable algorithm that allows us to efficiently compute gradients for even the most massive neural networks.
    
    \item[How do we choose the learning rate $\alpha$?] This is a critical hyperparameter. We'll discuss learning rate schedules and more advanced optimization algorithms.
    
    \item[How can we improve upon SGD?] We will explore modern optimizers like \textbf{SGD with Momentum} and \textbf{Adam}, which try to accelerate convergence and smooth out the noisy path of SGD.
\end{description}



\end{document}


% \documentclass[11pt]{article}
% \usepackage{amsmath, amssymb, amsfonts}
% \usepackage{geometry}
% \usepackage{enumerate}
% \usepackage{graphicx}
% \geometry{margin=1in}

% \title{Stochastic Gradient Descent (SGD) for Deep Learning\\
% \large A Blackboard Teaching Guide}
% \author{Based on Gilbert Strang's Linear Algebra and Learning from Data, Section VI.5}
% \date{}

% \begin{document}
% \maketitle

% \section{The Fundamental Problem Setup}

% \subsection{The Optimization Framework}
% In machine learning, we repeatedly encounter optimization problems of the form of a finite sum:
% \begin{equation}
% \min_{x} L(x) = \min_{x} \frac{1}{N} \sum_{i=1}^{N} \ell_i(x)
% \end{equation}

% Where:
% \begin{itemize}
% \item $x \in \mathbb{R}^d$ are the model parameters (weights in neural networks)
% \item $N$ is the number of training samples (often millions)
% \item $d$ is the parameter dimension (can be millions/billions)
% \item $\ell_i(x)$ is the loss on the $i$-th training sample
% \end{itemize}
% \textbf{Key Challenge:} Both $N$ and $d$ can be enormous in modern deep learning!


% Classical Examples:
% \begin{itemize}
%     \item \textbf{Least Squares:} $L(x) = \frac{1}{2}\|Ax - b\|^2 = \frac{1}{2N}\sum_{i=1}^{N}(a_i^T x - b_i)^2$
%     \item \textbf{Deep Neural Networks:} $L(x) = \frac{1}{N}\sum_{i=1}^{N} \ell(\text{DNN}(a_i; x), y_i)$
% \end{itemize}
% Where $\text{DNN}(a_i; x)$ represents the neural network output for input $a_i$ with parameters $x$.


% Note that $\ell_i(x)$  captures the entire NN architecture though a particular instance of the finite sum. Most of the problem in ML and statistics (MLE)when written down as an optimization problem looks like that finite sum problem. Hence we need specialized algorithm for such finite sum problems. 


% \section{From Gradient Descent to SGD}

% \subsection{Standard Gradient Descent}
% The classical gradient descent iteration is:
% \begin{equation}
% x_{k+1} = x_k - \alpha_k \nabla L(x_k)
% \end{equation}

% \textbf{Expanding the gradient:}
% \begin{equation}
% \nabla L(x_k) = \frac{1}{N} \sum_{i=1}^{N} \nabla \ell_i(x_k)
% \end{equation}

% \textbf{The Computational Bottleneck:} Computing $\nabla L(x_k)$ for only one step of gradient descent requires:
% \begin{itemize}
% \item Processing all $N$ training samples
% \item Computing $N$ individual gradients $\nabla \ell_i(x_k)$
% \item For large $N$ (millions), this can take hours per iteration!
% \end{itemize}

% \subsection{The SGD Solution}
% \textbf{Core Idea:} Instead of using the full gradient, use a single randomly chosen sample:

% \begin{equation}
% x_{k+1} = x_k - \alpha_k \nabla \ell_{i_k}(x_k)
% \end{equation}

% Where $i_k$ is chosen uniformly at random from $\{1, 2, \ldots, N\}$.

% \textbf{Computational Gain:} Each iteration is approximately $N$ times faster!

% \section{Why Does SGD Work? Mathematical Foundation}

% \subsection{The Unbiased Property}
% \textbf{Key Mathematical Insight:} The stochastic gradient is an unbiased estimator:

% \begin{equation}
% \mathbb{E}_{i_k}[\nabla \ell_{i_k}(x_k)] = \frac{1}{N}\sum_{i=1}^{N} \nabla \ell_i(x_k) = \nabla L(x_k)
% \end{equation}

% \textbf{Proof:} 
% \begin{align}
% \mathbb{E}_{i_k}[\nabla \ell_{i_k}(x_k)] &= \sum_{i=1}^{N} P(i_k = i) \nabla \ell_i(x_k)\\
% &= \sum_{i=1}^{N} \frac{1}{N} \nabla \ell_i(x_k) = \nabla L(x_k)
% \end{align}

% \textbf{Intuition:} On average, SGD steps point in the same direction as true gradient descent!

% \subsection{The Variance Trade-off}
% While unbiased, the stochastic gradient has higher variance:
% \begin{equation}
% \text{Var}[\nabla \ell_{i_k}(x_k)] > 0
% \end{equation}

% This variance causes the characteristic ``noisy" behavior of SGD.

% \section{Geometric Intuition: The 1D Case}

% \subsection{A Simple Quadratic Example}
% Consider the 1D problem: $L(x) = \frac{1}{N}\sum_{i=1}^{N} (a_i x - b_i)^2$

% \textbf{Individual minima:} Each $\ell_i(x) = (a_i x - b_i)^2$ is minimized at $x_i^* = \frac{b_i}{a_i}$

% \textbf{Global minimum:} $x^* = \frac{\sum a_i b_i}{\sum a_i^2}$ lies in the interval $[\min_i \frac{b_i}{a_i}, \max_i \frac{b_i}{a_i}]$

% \subsection{Two Regions of Behavior}

% \textbf{Far-Out Zone:} When $x_k$ is outside $[\min_i \frac{b_i}{a_i}, \max_i \frac{b_i}{a_i}]$
% \begin{itemize}
% \item Most stochastic gradients $\nabla \ell_{i_k}(x_k) = 2a_i(a_i x_k - b_i)$ have the same sign as $\nabla L(x_k)$
% \item SGD makes consistent progress toward the solution
% \item \textbf{Result:} Rapid initial convergence
% \end{itemize}

% \textbf{Region of Confusion:} When $x_k$ is inside the interval
% \begin{itemize}
% \item Different stochastic gradients can have different signs
% \item SGD oscillates around the true minimum
% \item \textbf{Result:} Noisy behavior near convergence
% \end{itemize}
% We may like the regions of confusion in ML for gaining robustness..

% In other words, we provide a randomized estimate of the gradient! But the key is in expectation (over that randomness), we have an unbiased estimate of gradient. So average is OK, but we need to discuss the variance too! A lot of research has been to still have an unbiased estimate but with less variance.


% The proof that SGD is well-behaved in both convex and nonconvex (in DNN case).


% \section{Practical SGD: Minibatches}
% Allows for parallel computation in GPU!
% \subsection{Minibatch SGD}
% Instead of single samples, use small batches:
% \begin{equation}
% x_{k+1} = x_k - \alpha_k \frac{1}{B} \sum_{j \in \text{Batch}_k} \nabla \ell_j(x_k)
% \end{equation}

% Where $B$ is the minibatch size (typically 32, 64, 128, 256).

% \subsection{Benefits of Minibatches}
% \begin{enumerate}
% \item \textbf{Variance Reduction:} $\text{Var}[\frac{1}{B}\sum_{j=1}^{B} \nabla \ell_j(x)] = \frac{1}{B}\text{Var}[\nabla \ell_j(x)]$
% \item \textbf{Computational Efficiency:} GPUs can process minibatches in parallel
% \item \textbf{Stable Convergence:} Less noisy than single-sample SGD
% \end{enumerate}

% \subsection{The Bias-Variance-Parallelism Trade-off}
% \begin{itemize}
% \item $B = 1$: Pure SGD (high variance, maximum stochasticity)
% \item $B = N$: Full batch GD (no variance, deterministic)
% \item very large minibatch is bad, it may seem interesting to maximize the GPU use, in DNN is like batch SGD, decreases noise so much that the region of confusion that sounds good but is bad for ML. Overfitting and sucking in test data!
% \item $B \in [32, 256]$: Sweet spot for deep learning
% \end{itemize}

% \section{Convergence Analysis}

% \subsection{Convergence in Expectation}
% For the stochastic gradient descent with step size $\alpha_k$:
% \begin{equation}
% \mathbb{E}[L(x_{k+1})] \leq \mathbb{E}[L(x_k)] - \frac{\alpha_k}{2}\mathbb{E}[\|\nabla L(x_k)\|^2] + \frac{\alpha_k^2}{2}\mathbb{E}[\|\nabla \ell_{i_k}(x_k)\|^2]
% \end{equation}

% \textbf{Key insight:} Progress depends on the balance between gradient magnitude and noise level.
% Because of its nature (only using a random sample or subset of sample instead of all data), sensitive to step size. 
% \subsection{Learning Rate Considerations}

% Choosing $\alpha$, is a hot research topic. We do not necessarily find an exact optimum in the training data! Sometimes we need a quick and dirty progress there!

% \textbf{Too large $\alpha_k$:} Amplifies noise, causes divergence

% \textbf{Too small $\alpha_k$:} Slow convergence

% \textbf{Common schedules:}
% \begin{itemize}
% \item Constant: $\alpha_k = \alpha$
% \item Decreasing: $\alpha_k = \frac{\alpha_0}{\sqrt{k}}$
% \item Step decay: $\alpha_k = \alpha_0 \gamma^{\lfloor k/s \rfloor}$
% \end{itemize}

% \section{Advanced Topics}
% A lot of research has been to still have an unbiased estimate but with less variance.

% The proof that SGD is well-behaved in both convex and nonconvex (in DNN case).


% \subsection{Two variance of Sampling Strategies}
% \textbf{With Replacement:} Sample uniformly at random each step
% \begin{itemize}
% \item Maintains independence between steps
% \item Easier theoretical analysis
% \item Less common in practice
% \end{itemize}

% \textbf{Without Replacement:} Shuffle dataset, iterate through
% \begin{itemize}
% \item Standard in practice (one ``epoch")
% \item Better data efficiency
% \item Harder theoretical analysis!!
% \end{itemize}

% \subsection{The Generalization Mystery}
% \textbf{Empirical Observation:} SGD often generalizes better than full batch methods!

% \textbf{Possible explanations:}
% \begin{itemize}
% \item Noise acts as implicit regularization
% \item SGD finds ``flatter" minima that generalize better
% \item Early stopping prevents overfitting
% \end{itemize}



% \subsection{Modern Variants}
% \textbf{SGD with Momentum:}
% \begin{align}
% v_{k+1} &= \beta v_k + \nabla \ell_{i_k}(x_k)\\
% x_{k+1} &= x_k - \alpha_k v_{k+1}
% \end{align}

% \textbf{Adam:} Combines momentum with adaptive learning rates

% \section{Summary}

% \textbf{SGD Key Points:}
% \begin{enumerate}
% \item Replaces expensive full gradients with cheap stochastic estimates
% \item Works because stochastic gradients are unbiased
% \item Exhibits fast initial progress, then noisy convergence
% \item Minibatches balance computational efficiency and variance reduction
% \item Noise may help generalization in deep learning
% \end{enumerate}

% \textbf{The Bottom Line:} SGD trades statistical efficiency for computational efficiency, making large-scale machine learning feasible.




% \subsection{How to compute stochastic gradients in Deep Learning?}

% We did not discuss how to compute  $\nabla \ell_i(x)$ in neural networks which will be the topic of next leacture: backpropagation:
% \begin{equation}
% \nabla \ell_i(x) = \nabla_x \ell(\text{DNN}(a_i; x), y_i)
% \end{equation}
% \end{document}